\section{Additional validation and provisional analyses}
The new figures expose available numerical results rather than replacing unvalidated locations with blank panels. Independent generation does not establish their scientific accuracy. The frozen calibration cohort and original 3D grid are unchanged; the retained original grid hash and conventions are documented above.

\paragraph{Precision and convergence diagnosis.}
The original 32/64-bit initial model, station file, observations and control files are byte-identical. Both prelocation runs stop at 14 iterations. Their maximum depth difference is initially 0.02 km and grows to 0.89 km through coupled updates. Common-start tests give identical printed round-1 locations and zero printed round-2 depth difference; the latter horizontal difference is 11.1 m. The small final update can result from four-backup stagnation and is not proof of an optimum. Backup decisions diverge in round 2; subsequent paths stop either at iteration limits or after four backups. Instrumented, otherwise unchanged generic solvers reproduce the original output bytes. The damped prelocation normal-matrix condition numbers range from 2372 to 3313; this test does not support a catastrophic linear-solve failure. Nonlinear paths, ray changes and two-decimal restart serialization amplify initial differences. Exact arithmetic attribution of every ray transition remains untested.

Repeated fixed-model prelocation, up to ten 14-iteration passes in each precision, still fails the unchanged 100 m stability gate. The 64-bit last pass moves an event 282 m horizontally and 0.24 km vertically. More iterations alone do not resolve this oscillation. No locations from this failed correction are substituted. Every fifth calibration event was withheld: 30 events train the model, 8 events (126 phases) test it. Unweighted held-out phase RMS is 0.2177 s for the starting profile and 0.2198 s for the refitted profile. A smaller solver event statistic is a different quantity and cannot reverse this failed prediction check. Manual phase/association review and absolute-location validation remain necessary. The own fitted 1--3 km layer has Vp/Vs=1.17. Passing positive-velocity ordering alone does not establish a plausible or well-constrained elastic structure; this low ratio warrants scientific review.

\paragraph{Relocation and conditional uncertainty.}
Across 124 common events, rounds 1 and 2 differ by a median 113.3 m and maximum 384.0 m horizontally; median absolute depth change is 0.095 km and maximum 0.382 km. On the same 399 waveform constraints, RMS improves from 0.0414 to 0.0400 s. The same 4395 picked-phase constraints improve from 0.1084 to 0.1020 s. Fixed grid, positive travel times, all-cluster residual coverage and recomputed station-gap checks pass. These are internal checks, not absolute accuracy. 34 events have no retained waveform correlations. No event passes S9's $NCCP+NCCS>100$ filter.

SVD was attempted on every connected component with at least four events in four explicit quadrants; 90 estimates were returned. Their 2-sigma depth errors reach 4.002 km, whereas the publication describes most coordinate errors near 10 m. The regional masks, sample sizes and starting locations differ, so this is a consequential failure to recover comparable precision, not a contradiction of the original estimates. S7 retains the large uncertainties. S12 uses 118 recorded component traces joined to the same relocated cohort. Sector masks and spectral normalization are explicit assumptions; no published depth bands select our events.

\paragraph{Magnitude and source mechanisms.}
The relocated-subset magnitude stage measured 5929 channels, with 1626 passing 5 dB, and inverted 884 relative-amplitude equations to obtain 65 magnitudes. Routine duration magnitudes enter only as regularization after independent detection and location. All accepted components have full numerical rank; formal errors are conditional on the response, location, weighting and unit regularization assumptions. Maximum curvature gives $M_c\simeq 0.0$ and the conditional likelihood estimate is $b=0.866\pm 0.110$ (41 above $M_c$). This selected, discontinuously observed sample cannot establish the published full-catalog $M_c=0.4$ or $b=0.97$.

SKHASH 1.1.5 now executes on 76 automated measurements from 8 events; 35 pass an explicit sign/SNR screen. Seven events fail the minimum-eight-polarity rule, and event 641 fails the unchanged azimuthal-gap rule (186.158 versus 90 degrees). No mechanism is returned. These are unreviewed candidate polarities, not Tan's manual measurements; no supplied mechanism or S/P value is used. Independent gCAP/FK processing, reviewed S/P measurements and event-specific waveform settings remain incomplete. The generic software is isolated in \path{.venv-skhash} because its NumPy requirement differs from the main analysis environment; the main NumPy 2.4.6 was restored. Dependency caches remain under the existing user pip cache; scientific products remain inside this repository.

\paragraph{Commands and evidence.}\begingroup\small\raggedright
\path{scripts/diagnose_grid_paths.py}, \path{scripts/settle_grid_prelocation.py}, \path{scripts/grid_uncertainty.py}, \path{scripts/downstream_checks.py --dataset grid}, \path{scripts/grid_focal_attempt.py} and \path{scripts/grid_panel_data.py} generate the new saved products. The single procedure \path{.venv/bin/python scripts/rebuild_provisional_workflow.py} executes these stages and compiles both documents. \path{docs/PROVISIONAL_CONTINUATION.md} records the full command and input chain, including acquisition scope. Procedure completion is separate from scientific acceptance.
\par\endgroup
